\subsection{Reinforcement Learning Implementation}

This subsection explains the process of training the Q-learning controller, embedding
it into a compact format suitable for use in the embedded platform and integrating it
into the actual alignment procedure. The goal of this design was to make sure that the
policy itself would be light enough to run on the ESP32-S3 chip, yet sufficiently rich
to cope with the RSSI noisiness and varying starting orientation.

\subsubsection{State, Action, and Reward Representation}

In order to implement the reinforcement learning algorithm, we need to define the
relevant state-space and actions:
\begin{equation}
s_k = \left( i_{\text{pan}},\; i_{\text{tilt}},\; \text{trend}_k \right)
\end{equation}
Here, $i_{\text{pan}}$ and $i_{\text{tilt}}$ are discretizations of pan and tilt indices,
respectively, while the third component stands for the short-term RSSI trend. Such state
definition allowed to leverage the information about orientation and current RSSI value
without creating a cumbersome state vector.

We defined the following set of actions for both training and inference stages:
\begin{equation}
\mathcal{A} = \{(+\Delta_p,0), (-\Delta_p,0), (0,+\Delta_t), (0,-\Delta_t), (0,0)\}
\end{equation}
where $\Delta_p$ and $\Delta_t$ denote the corresponding pan and tilt increments. The
hold action was necessary for ensuring that the alignment stops moving when it enters
high RSSI area.

The reward is given by the following expression:
\begin{equation}
r_k = \Delta \text{RSSI}_k - \lambda \cdot \mathbb{I}_{\text{move}}
\end{equation}
This reward shaping improved convergence smoothness and reduced oscillatory behavior in
both simulation and practical inference.

\subsubsection{Training Setup}

The training was performed offline in the Python environment by repeating episodes from random starting positions. The hyperparameters were configured to enable the agent to converge without any overfitting toward one start position.

\begin{table}[H]
\centering
\caption{Q-Learning Training Configuration}
\label{tab:q_learning_parameters}
\renewcommand{\arraystretch}{1.4}
\begin{tabular}{|p{4.8cm}|c|p{5.0cm}|}
\hline
\textbf{Parameter} & \textbf{Value} & \textbf{Purpose} \\
\hline
Number of training episodes & 2500 & Provides repeated exposure to diverse initial orientations \\
\hline
Maximum steps per episode & 60 & Prevents long unproductive trajectories \\
\hline
Learning rate $\alpha$ & 0.18 & Balances adaptation speed and numerical stability \\
\hline
Discount factor $\gamma$ & 0.92 & Encourages actions with long-term RSSI improvement \\
\hline
Initial exploration rate $\epsilon$ & 1.0 & Supports broad exploration at the start of training \\
\hline
Final exploration rate & 0.05 & Preserves limited exploration near convergence \\
\hline
Exploration decay & 0.995 per episode & Gradually shifts learning toward exploitation \\
\hline
Movement penalty $\lambda$ & 0.35 & Discourages unnecessary actuator motion \\
\hline
\end{tabular}
\end{table}

\subsubsection{Storing the Q-Table and Embedding It into the System}

Upon completion of the learning process, the Q-table obtained through Python code was transformed to a static array structure to facilitate embedding it into firmware. The aim was to prevent dynamic file parsing at runtime by the microcontroller.

\begin{table}[H]
\centering
\caption{Embedded RL Deployment Summary}
\label{tab:rl_deployment_summary}
\renewcommand{\arraystretch}{1.4}
\begin{tabular}{|p{3.6cm}|c|p{5.3cm}|}
\hline
\textbf{Item} & \textbf{Value} & \textbf{Remarks} \\
\hline
Pan states & 19 & Practical deployment range of $0^{\circ}$--$180^{\circ}$ at $10^{\circ}$ resolution \\
\hline
Tilt states & 11 & Practical deployment range of $25^{\circ}$--$75^{\circ}$ at $5^{\circ}$ resolution \\
\hline
RSSI trend states & 3 & Increase, stable, and decrease \\
\hline
Actions per state & 5 & Left, right, up, down, hold \\
\hline
Total Q-values & 3135 & Compact enough for direct static allocation \\
\hline
Stored numeric type & 32-bit float & Preserves training precision with low memory overhead \\
\hline
Estimated memory footprint & 12.3~kB & Well within ESP32-S3 memory constraints \\
\hline
Inference policy & Greedy $\arg\max Q(s,a)$ & No online learning during hardware testing \\
\hline
\end{tabular}
\end{table}

Inference was the sole method employed by the firmware in practice. On-line
reinforcement learning based updates to Q-values were not performed to avoid the potential
instability stemming from noise in packet-level metrics, to minimize computation time, and to
ensure repeatability across trials.

\subsubsection{Operational Sequence During Practical Trials}

At runtime, the controller based on RL was run through a two-step process.
First, the antenna was positioned into a valid initial position and RSSI samples were collected. Second,
the controller would continuously encode the input, choose the highest valued action according to the
contents of the Q-table, and issue the corresponding commands to the motors until convergence.

Convergence was achieved in either one of these cases:
\begin{itemize}
\item If the absolute RSSI gain was smaller than 0.8 dB for 4 consecutive iterations, or 
\item When the maximum number of steps had been reached by the algorithm.
\end{itemize}
